\documentclass[11pt]{article}
\usepackage{cmap}
\usepackage[margin=1.05in]{geometry}
\usepackage{amsmath,amssymb,amsthm}
\usepackage[T1]{fontenc}
\usepackage{lmodern}
\usepackage{iftex}
\ifPDFTeX
\input{glyphtounicode}
\pdfgentounicode=1
\fi
\usepackage{xcolor}
\usepackage{booktabs}
\usepackage{microtype}
\usepackage{hyperref}
\definecolor{linkblue}{rgb}{0.08,0.19,0.52}
\hypersetup{colorlinks=true,linkcolor=linkblue,citecolor=linkblue,
  urlcolor=linkblue,pdftitle={Endpoint Parity Loss in a Bessel Wronskian},
  pdfauthor={Lluis Eriksson}}

\newtheorem{theorem}{Theorem}
\newtheorem{lemma}[theorem]{Lemma}
\newtheorem{proposition}[theorem]{Proposition}
\newtheorem{corollary}[theorem]{Corollary}
\theoremstyle{definition}
\newtheorem{conjecture}[theorem]{Conjecture}
\theoremstyle{remark}
\newtheorem{remark}[theorem]{Remark}

\newcommand{\FA}{F_{\!A}}
\newcommand{\FB}{F_{\!B}}
\newcommand{\W}{\mathcal W}
\newcommand{\dd}{\,d}
\newcommand{\I}{I}

\title{Endpoint Parity Loss in a Bessel Wronskian:\\
\large an exact obstruction to kernel-and-anchor proofs of global ratio monotonicity}
\author{Lluis Eriksson}
\date{August 2026}

\begin{document}
\maketitle

\begin{abstract}
For $\beta>0$ write $I_m=I_m(\beta)$ and define
\[
 a_m=I_m^2\bigl((m-1)I_{m-1}^2+(m+1)I_{m+1}^2\bigr),
 \qquad b_m=mI_m^4,
\]
with sine series $\FA(t)=\sum_{m\geq1}a_m\sin(mt)$ and
$\FB(t)=\sum_{m\geq1}b_m\sin(mt)$.  The associated Wronskian
$\W=2(\FA'\FB-\FA\FB')$ is negative exactly when $\FA/\FB$ is
decreasing.  We isolate what can and cannot be proved from two natural
inputs: the Neumann convolution kernel
$I_0(2\beta\sin(\phi/2))$ and the small-coupling anchor
$(-4096/\beta^{10})\W\to4\sin^3t$.

First, the convolution does prove the denominator statement
$\FB(t)>0$ for $0<t<\pi$.  Second, the endpoint is governed by two
alternating quantities, $c_3$ and $B_\pi$, through the exact law
$\W(\pi-d)=-4c_3B_\pi d^3+O(d^5)$.  We prove $B_\pi>0$, derive integral
representations, and establish the logarithmic cancellation law
\[
 \lim_{\beta\to\infty}-\frac{1}{\beta}\log
 \frac{|c_3B_\pi|}{c_3^{\rm abs}B_\pi^{\rm abs}}
 =8-4\sqrt2.
\]
Thus a positive-term majorant loses an exponentially decisive factor.
Finally, we give a smooth one-parameter perturbation theorem: one may keep
$\FB$ and its kernel unchanged, preserve positivity and strict coefficient-ratio
ordering, and preserve every jet at $\beta=0$, while choosing either sign of
the endpoint cubic coefficient.  Consequently those structural data, even
taken together, do not imply global Wronskian negativity.  This is a no-go
theorem for a proof architecture, not a counterexample to the Bessel
conjecture.  A short high-precision kill-test accompanies the paper; no
computer-assisted inequality is used in the proofs.
\end{abstract}

\section{The question and the status of the result}

The following ratio-monotonicity statement arises after an exact
Wronskian reduction of a surface-positivity problem.

\begin{conjecture}[Bessel Wronskian conjecture]\label{conj:main}
For every $\beta>0$,
\[
 \FB(t)>0,
 \qquad
 \left(\frac{\FA}{\FB}\right)'(t)<0
 \quad (0<t<\pi).
\]
Equivalently, $\FB>0$ and $\W<0$ on $(0,\pi)$.
\end{conjecture}

The first assertion will be proved exactly below.  The second assertion is
\emph{not} proved here.  Our main result instead separates a valid part of
the proposed kernel route from a fatal loss at $t=\pi$.  The distinction is
important: the perturbations in Theorem~\ref{thm:nogo} do not obey all
modified-Bessel recurrences, and hence do not refute
Conjecture~\ref{conj:main}.

The Wronskian identity itself is elementary.  If
$F_x(t)=\sum_{m\ge1}x_m\sin(mt)$, then
\begin{equation}\label{eq:wronskian}
 2(F_a'F_b-F_aF_b')=
 \sum_{m<n}(a_mb_n-a_nb_m)
 \bigl[(m-n)\sin((m+n)t)+(m+n)\sin((n-m)t)\bigr].
\end{equation}
It follows by expanding the products and antisymmetrizing.  Throughout,
all Bessel series and their differentiated series converge locally
uniformly; the rearrangements below are therefore legitimate.

\section{The convolution kernel closes \texorpdfstring{$\FB>0$}{FB > 0}}

The Neumann addition theorem gives the two parity forms
\begin{align}
 K_+(\phi)&=I_0\!\left(2\beta\cos\frac{\phi}{2}\right)
   =\sum_{m\in\mathbb Z}I_m(\beta)^2e^{im\phi},\label{eq:kplus}\\
 K_-(\phi)&=I_0\!\left(2\beta\sin\frac{\phi}{2}\right)
   =\sum_{m\in\mathbb Z}(-1)^mI_m(\beta)^2e^{im\phi}.
   \label{eq:kminus}
\end{align}
For periodic convolution normalized by $(f*g)(t)=(2\pi)^{-1}
\int_0^{2\pi}f(s)g(t-s)\,ds$, both squares give
\begin{equation}\label{eq:p4}
 P_4:=K_+*K_+=K_-*K_-
 =\sum_{m\in\mathbb Z}I_m(\beta)^4e^{imt},
 \qquad \FB=-\frac12P_4'.
\end{equation}

\begin{lemma}[strict circular convolution]\label{lem:circle}
If $q$ is a nonconstant, even, continuous function on the circle that is
strictly decreasing with geodesic distance from $0$, then $q*q$ is strictly
decreasing on $(0,\pi)$ and $(q*q)'<0$ there whenever $q$ is smooth.
\end{lemma}

\begin{proof}
Use the layer-cake representation of $q-\min q$ by centered circular arcs.
The convolution is a positive superposition of overlap lengths of two such
arcs whose centers are separated by $t$.  Each overlap length is
nonincreasing for $0<t<\pi$.  At any fixed $t$ there is an open set of pairs
of level arcs whose endpoints meet in the unsaturated regime; on this set
the derivative of the overlap is strictly negative.  Strict decrease of
$q$ makes this set have positive level measure.  Differentiation is allowed
for smooth $q$ and yields the pointwise derivative claim.
\end{proof}

\begin{theorem}\label{thm:fb}
For every $\beta>0$, $\FB(t)>0$ for $0<t<\pi$.
\end{theorem}

\begin{proof}
On $[0,\pi]$, $K_+(t)=I_0(2\beta\cos(t/2))$ is strictly decreasing,
because $I_0'(x)=I_1(x)>0$ for $x>0$.  Apply
Lemma~\ref{lem:circle} to \eqref{eq:p4}: $P_4'(t)<0$, and hence
$\FB(t)=-P_4'(t)/2>0$.
\end{proof}

This proof uses the kernel at full strength and contains no alternating
majorization.  The obstruction begins with the numerator and becomes sharp
at the opposite endpoint.

\section{The exact small-coupling anchor}

The fixed-order Taylor expansions of $I_m$ give
\begin{align*}
 a_1&=\frac{\beta^6}{128}+O(\beta^8),&
 b_1&=\frac{\beta^4}{16}+O(\beta^6),\\
 a_2&=\frac{\beta^6}{256}+O(\beta^8),&
 b_2&=\frac{\beta^8}{2048}+O(\beta^{10}).
\end{align*}
All remaining modes enter later.  Substitution in the Wronskian yields,
uniformly for $t\in[0,\pi]$,
\begin{equation}\label{eq:anchor}
 \W(t;\beta)=-\frac{\beta^{10}}{1024}\sin^3t+O(\beta^{12}),
 \qquad
 -\frac{4096}{\beta^{10}}\W(t;\beta)\longrightarrow4\sin^3t.
\end{equation}
Thus the parabolic anchor has the desired sign for every fixed interior
$t$ when $\beta$ is small.  Uniformity of the absolute error in
\eqref{eq:anchor}, however, does not supply a relative estimate near
$t=\pi$, where the leading term has a triple zero.

\section{The endpoint coefficient}

Put $d=\pi-t$.  The linear coefficient of $\FA(\pi-d)$ vanishes by an
edgewise telescoping:
\begin{equation}\label{eq:telescoping}
 \sum_{m\ge1}(-1)^{m+1}ma_m=0.
\end{equation}
Indeed, the two appearances of $I_m^2I_{m+1}^2$ have coefficients
\[
 (-1)^{m+1}m(m+1)
 \quad\hbox{and}\quad
 (-1)^{m+2}(m+1)m,
\]
which cancel.  Define
\begin{align}
 c_3(\beta)&=\frac{1}{6}\sum_{m\ge1}(-1)^{m+1}
 m(m+1)(2m+1)I_m^2I_{m+1}^2,\label{eq:c3}\\
 B_\pi(\beta)&=\sum_{m\ge1}(-1)^{m+1}m^2I_m^4.\label{eq:bpi}
\end{align}
Taylor expansion, followed by the same edge collection used in
\eqref{eq:telescoping}, gives
\begin{equation}\label{eq:endexp}
 \FA(\pi-d)=c_3d^3+O(d^5),\qquad
 \FB(\pi-d)=B_\pi d+O(d^3),
\end{equation}
and consequently
\begin{equation}\label{eq:wend}
 \W(\pi-d)=-4c_3B_\pi d^3+O(d^5),
 \qquad
 \lim_{t\uparrow\pi}\frac{-\W(t)}{\sin^3t}=4c_3B_\pi.
\end{equation}

The denominator coefficient has a positive integral.  From
$P_4''(\pi)=2B_\pi$, differentiating the convolution and integrating by
parts once gives
\begin{equation}\label{eq:bpi-integral}
 B_\pi=\frac{\beta^2}{\pi}\int_0^{\pi/2}
 \sin u\cos u\,I_1(2\beta\cos u)I_1(2\beta\sin u)\,du>0.
\end{equation}
Therefore the desired endpoint sign is exactly the nontrivial statement
$c_3(\beta)>0$.

For later use, set $f(x)=I_1(x)/2$, $y=2\beta\sin u$ and
$z=2\beta\cos u$.  Applying \eqref{eq:kminus}, differentiating the
correlation, and integrating by parts gives the half-range identity
\begin{equation}\label{eq:c3-integral}
 c_3=\frac{1}{6\pi}\int_0^{\pi/2}J(u,\beta)\,du,
\end{equation}
where
\begin{equation}\label{eq:J}
 J=-z^2f'(y)f'(z)+zy^2f'(y)f''(z)+yf(y)f'(z).
\end{equation}
The integrand is not pointwise positive.  This is not a numerical
accident: for every sufficiently large $\beta$, its negative boundary
region and positive saddle region are both nonempty.  Hence
\eqref{eq:c3-integral} is a signed cancellation identity rather than a
positive-kernel proof of $c_3>0$.

\section{The exponential parity loss}

Introduce the positive envelopes
\begin{align}
 c_3^{\rm abs}&=\frac{1}{6}\sum_{m\ge1}
 m(m+1)(2m+1)I_m^2I_{m+1}^2,\label{eq:cabs}\\
 B_\pi^{\rm abs}&=\sum_{m\ge1}m^2I_m^4.\label{eq:babs}
\end{align}

\begin{theorem}[endpoint cancellation exponent]\label{thm:exponent}
As $\beta\to\infty$,
\begin{align}
 c_3(\beta)&=\frac{\beta^{3/2}e^{2\sqrt2\beta}}
  {24\,2^{1/4}\pi^{3/2}}\bigl(1+O(\beta^{-1})\bigr),
  \label{eq:c3asymp}\\
 B_\pi(\beta)&=\frac{\beta^{1/2}e^{2\sqrt2\beta}}
  {4\,2^{3/4}\pi^{3/2}}\bigl(1+O(\beta^{-1})\bigr).
  \label{eq:bpiasymp}
\end{align}
Moreover
\begin{equation}\label{eq:absrate}
 \lim_{\beta\to\infty}\frac{1}{\beta}\log c_3^{\rm abs}
 =\lim_{\beta\to\infty}\frac{1}{\beta}\log B_\pi^{\rm abs}=4,
\end{equation}
and therefore
\begin{equation}\label{eq:loss}
 \boxed{\displaystyle
 \lim_{\beta\to\infty}-\frac{1}{\beta}\log
 \frac{|c_3B_\pi|}{c_3^{\rm abs}B_\pi^{\rm abs}}
 =8-4\sqrt2.}
\end{equation}
\end{theorem}

\begin{proof}
We give the details needed to identify both the sign and the exponential
scale.  In \eqref{eq:c3-integral}, the phase is
$y+z=2\beta(\sin u+\cos u)$.  It has a unique nondegenerate maximum at
$u_0=\pi/4$.  Write $x=\sqrt2\beta$ and $u=u_0+v$.  Uniformly for
$|v|\leq\beta^{-2/5}$, the fixed-order large-argument expansion and its
differentiated forms give
\[
 f^{(j)}(y)=\frac{e^y}{2\sqrt{2\pi y}}
 (1+O(\beta^{-1}+|v|)),\qquad j=0,1,2,
\]
and the analogous statement at $z$.  The leading term in $J$ is
$zy^2f(y)f(z)$; at the saddle it equals
\[
 \frac{x^2e^{2x}}{8\pi}\bigl(1+O(\beta^{-1}+|v|)\bigr).
\]
Also $y+z=2x-xv^2+O(\beta v^4)$.  Laplace's method in
\eqref{eq:c3-integral} now gives
\[
 c_3=\frac{1}{6\pi}\frac{x^2e^{2x}}{8\pi}
 \sqrt{\frac{\pi}{x}}\,(1+O(\beta^{-1})),
\]
which is \eqref{eq:c3asymp}.  Outside the displayed central window,
$\sin u+\cos u$ has a uniform quadratic deficit until a fixed
neighborhood is exited and a uniform positive deficit thereafter; the
standard exponential Bessel bound makes that contribution smaller than
the stated error.  This also justifies the differentiated asymptotics.

The same saddle applied to \eqref{eq:bpi-integral} is simpler.  At $u_0$
the prefactor $\sin u\cos u$ equals $1/2$, and
$I_1(x)^2=e^{2x}(2\pi x)^{-1}(1+O(x^{-1}))$.  The Gaussian factor
$\sqrt{\pi/x}$ yields \eqref{eq:bpiasymp}.

For \eqref{eq:absrate}, any fixed positive summand gives the lower
exponential rate $4$, since $I_m(\beta)=e^\beta
(2\pi\beta)^{-1/2}(1+O(\beta^{-1}))$ for fixed $m$.  For the upper bound,
the generating function
\[
 \sum_{m\in\mathbb Z}I_m(\beta)z^m
 =\exp\!\left(\frac{\beta}{2}(z+z^{-1})\right)
\]
and the three-term recurrence control every fixed absolute polynomial
moment.  Positivity and repeated product bounds dominate
\eqref{eq:cabs} and \eqref{eq:babs} by a polynomial in $\beta$ times
$e^{4\beta}$.  Taking logarithms proves \eqref{eq:absrate}.  Combining
the two signed rates $2\sqrt2$ with the two absolute rates $4$ proves
\eqref{eq:loss}.
\end{proof}

The rate $8-4\sqrt2=2.3431457505\ldots$ is the kill-test for any
positive-term or absolute-value barrier.  A barrier whose slack decays only
algebraically cannot resolve the endpoint sign at large $\beta$.

\section{A structural no-go theorem}

We now show that the loss is logical, not merely a weakness of a particular
estimate.  Let $p=(2,1,0,\ldots)$.  It satisfies
\begin{equation}\label{eq:p}
 \sum_{m\ge1}(-1)^{m+1}mp_m=0,
 \qquad
 -\frac{1}{6}\sum_{m\ge1}(-1)^{m+1}m^3p_m=1.
\end{equation}

\begin{theorem}[structural non-identifiability]
\label{thm:nogo}
Let $(A_m(\beta),B_m(\beta))=(a_m(\beta),b_m(\beta))$ be the Bessel
arrays above, and suppose the known strict ordering
$A_m/B_m<A_{m+1}/B_{m+1}$ is included among the structural inputs.
There are two families $\widetilde A^+$ and $\widetilde A^-$ such that:
\begin{enumerate}
\item $B$ is unchanged; hence the convolution representation and
      $\FB>0$ are unchanged;
\item every $\widetilde A_m^\pm$ is positive and
      $\widetilde A_m^\pm/B_m$ remains strictly increasing in $m$ for
      every $\beta>0$;
\item $\widetilde A^\pm-A$ is flat at $\beta=0$, so the two families
      have the same complete small-$\beta$ asymptotic expansion and in
      particular the same anchor \eqref{eq:anchor};
\item at one common $\beta=\beta_*$, the endpoint cubic coefficients
      have opposite signs.  Thus their Wronskians have opposite signs
      in a punctured left neighborhood of $t=\pi$.
\end{enumerate}
Consequently items 1--3 cannot imply global Wronskian negativity.
\end{theorem}

\begin{proof}
For a scalar function $h$, put
\[
 \widetilde A_m=A_m+h(\beta)p_m,
 \qquad \widetilde B_m=B_m.
\]
By \eqref{eq:p}, the endpoint linear cancellation is preserved and the
cubic coefficient becomes $\widetilde c_3=c_3+h$.

Choose a large center $\beta_*=R$.  Theorem~\ref{thm:exponent} shows
$|c_3(R)|=O(R^{3/2}e^{2\sqrt2R})$.  In contrast, the fixed-mode
coefficients $A_1,A_2,B_1,B_2$ have exponential scale $e^{4R}$ times
an inverse power of $R$.  Direct insertion of the fixed-order Bessel
expansions gives, for fixed $m$,
\[
 \frac{A_{m+1}}{B_{m+1}}-\frac{A_m}{B_m}
 =\frac{4(2m+1)}{R^2}+O(R^{-3}).
\]
Multiplying these gaps by the relevant $B_m$ shows that the admissible
size of a perturbation in modes $1,2$ has exponential scale $e^{4R}$
times an inverse power of $R$.  Therefore, for all sufficiently large
$R$, there is an $\eta_R>|c_3(R)|$ such that $|h(R)|<\eta_R$ preserves
positivity and both affected ratio inequalities.

Fix $0<\varepsilon<\eta_R-|c_3(R)|$ and, for $\sigma>0$, set
\[
 \psi_{R,\sigma}(\beta)=
 \exp\!\left(-\frac{1}{\beta^2}-\frac{(\beta-R)^2}{\sigma^2}
             +\frac{1}{R^2}\right),
 \qquad
 h_\pm(\beta)=(-c_3(R)\pm\varepsilon)\psi_{R,\sigma}(\beta).
\]
Here $\psi(R)=1$; it is real analytic on $(0,\infty)$ and extends by
$0$ to a $C^\infty$ function flat at $0$.  At $R$ we have
$\widetilde c_3^\pm(R)=\pm\varepsilon$.

It remains only to choose $\sigma$.  In a neighborhood of $R$, strict
positivity and the two strict ratio gaps persist by continuity and the
choice of $\eta_R$.  On compact sets separated from that neighborhood,
$\psi_{R,\sigma}\to0$ uniformly as $\sigma\downarrow0$.  Near $0$ it is
smaller than every power of $\beta$, whereas the coefficients and their
strict gaps have finite algebraic orders.  Near infinity its Gaussian
decay dominates the fixed-mode coefficient scales and gap denominators.
Thus a sufficiently small $\sigma$ preserves positivity and ratio order
for all $\beta>0$.  Flatness gives item 3, and \eqref{eq:wend} gives item
4.
\end{proof}

\begin{remark}[exact scope]
The theorem rules out any argument whose hypotheses see only the $B$-kernel,
positivity of the coefficients, strict ordering of $A_m/B_m$, and the full
small-coupling jet.  A successful proof of Conjecture~\ref{conj:main} must
use additional Bessel structure that rejects the perturbation, for example
an exact signed identity, recurrence-compatible differential inequality,
or a cancellation-preserving certificate.
\end{remark}

\section{Reproducible kill-test}

The companion script
\texttt{scripts/wronskian\_endpoint\_kill\_test.py} is a diagnostic, not
an interval proof.  It uses 180-decimal arithmetic, truncates well beyond
the active Bessel modes, and checks the exact normalizations before
measuring the cancellation.  A representative run is:
\begin{center}
\begin{tabular}{@{}rrrr@{}}
\toprule
$\beta$ & endpoint ratio & $\log$(ratio) & next fitted rate\\
\midrule
8  & $3.86868117\,10^{-5}$  & $-10.1600118$ & $2.1068310$\\
16 & $1.85217583\,10^{-12}$ & $-27.0146600$ & $2.2301610$\\
32 & $5.90114568\,10^{-28}$ & $-62.6972361$ & $2.2878524$\\
64 & $9.45545116\,10^{-60}$ & $-135.9085142$ & ---\\
\bottomrule
\end{tabular}
\end{center}
Here ``endpoint ratio'' is
$|c_3B_\pi|/(c_3^{\rm abs}B_\pi^{\rm abs})$.  The fitted rates approach
$8-4\sqrt2$.  At $\beta=32$ and $d=10^{-4}$, the quotient of
$\W(\pi-d)/d^3$ by $-4c_3B_\pi$ is $1.0000009109773$.  The finite
perturbation that flips $c_3$ has relative size
$c_3/A_1=1.270694641\,10^{-12}$ and preserves positivity and the checked
ratio order.

The acceptance checks use explicit exceptions rather than Python
\texttt{assert}; the same run is required under ordinary Python and
\texttt{python -O}.  The option \texttt{--self-test-mutations} injects a
false predicate and requires the harness to reject it.  None of these
floating-point observations is used to establish
Theorems~\ref{thm:fb}, \ref{thm:exponent}, or \ref{thm:nogo}.

\section{Implications for proof design}

The two proposed legs now have sharply different outcomes.
\begin{enumerate}
\item The convolution kernel is strong enough to prove $\FB>0$ globally.
It should be retained as the denominator lemma.
\item The parabola $4\sin^3t$ is an exact small-$\beta$ anchor but is not a
global barrier by itself.  Any continuation in $(t,\beta)$ must carry an
endpoint variable normalized by $c_3B_\pi$, not merely absolute Fourier
moments.
\item The large-$\beta$ endpoint window requires signed resolution on the
scale $e^{-(8-4\sqrt2)\beta}$ relative to the naive positive envelope.
Polynomially accurate structural bounds necessarily miss it.
\item Coefficient-ratio order is useful but logically insufficient for
sine-series ratio monotonicity.  Theorem~\ref{thm:nogo} remains true even
after adjoining all jets at $\beta=0$.
\end{enumerate}

A plausible surviving architecture is therefore modular: use the exact
convolution for $\FB$; isolate $c_3>0$ by a cancellation-preserving signed
identity; prove a normalized endpoint window after factoring $\sin^3t$;
and bridge that window to a compact interior region.  The present paper
does not claim that this program has been completed.

\section{Priority and limitations}

The addition formulas \eqref{eq:kplus}--\eqref{eq:kminus} are classical
Neumann--Graf identities; standard sources are Watson's treatise and the
NIST Digital Library of Mathematical Functions~\cite{DLMFaddition,Watson}.
The large-argument expansions with differentiable error control are also
classical~\cite{DLMFasymptotic,Olver}.  Ratio criteria for ordinary power
series go back at least to Biernacki and Krzy\.z~\cite{BK}; they do not
directly apply to sine series.  Positivity and monotonicity phenomena for
Bessel sums have a substantial earlier literature, including
Askey--Steinig~\cite{AskeySteinig}.

The contribution claimed here is narrower: the exact endpoint reduction,
the cancellation exponent \eqref{eq:loss}, and the perturbative
non-identifiability theorem for the specified proof architecture.  We do
not claim novelty for the classical identities or Laplace method.  We also
do not claim a proof or disproof of Conjecture~\ref{conj:main}.  Related
programme manuscripts record the Wronskian reduction and a separate
certificate-based route to the global inequality~\cite{WRnote,Surface};
the present result is useful precisely because it explains why a much
shorter kernel-and-anchor proof cannot close without a new signed lemma.

\section*{Data and verification statement}

The source, diagnostic script, invocation commands, output transcript,
byte counts, and SHA-256 hashes are distributed with the manuscript.  The
paper proofs are ordinary analytic arguments.  No Lean build, interval
oracle, or sustained remote computation was used for this paper.  The
numerical kill-test is labelled diagnostic and is reproducible on a single
CPU process.

\begin{thebibliography}{99}
\small
\setlength{\itemsep}{2pt plus 0.2pt}

\bibitem{AskeySteinig}
R.~Askey and J.~Steinig,
Some positive trigonometric sums,
\emph{Transactions of the American Mathematical Society} \textbf{187}
(1974), 295--307.
\href{https://doi.org/10.1090/S0002-9947-1974-0338481-3}
{doi:10.1090/S0002-9947-1974-0338481-3}.

\bibitem{BK}
M.~Biernacki and J.~Krzy\.z,
On the monotonicity of certain functionals in the theory of analytic
functions,
\emph{Annales Universitatis Mariae Curie-Sk\l odowska, Sectio A}
\textbf{9} (1955), 5--23.
\href{https://bc.umcs.pl/dlibra/publication/33796/edition/30632}
{UMCS digital repository record}.

\bibitem{DLMFaddition}
NIST Digital Library of Mathematical Functions,
\S10.44(ii), addition theorems for modified Bessel functions,
\href{https://dlmf.nist.gov/10.44.ii}{dlmf.nist.gov/10.44.ii},
accessed 4 August 2026.

\bibitem{DLMFasymptotic}
NIST Digital Library of Mathematical Functions,
\S10.40, asymptotic expansions for large argument,
\href{https://dlmf.nist.gov/10.40}{dlmf.nist.gov/10.40},
accessed 4 August 2026.

\bibitem{Olver}
F.~W.~J. Olver,
\emph{Asymptotics and Special Functions},
A K Peters, Wellesley, MA, 1997.

\bibitem{Watson}
G.~N. Watson,
\emph{A Treatise on the Theory of Bessel Functions}, 2nd ed.,
Cambridge University Press, 1944.

\bibitem{WRnote}
L.~Eriksson,
A Wronskian identity for antisymmetrized sine sums,
programme manuscript, July 2026.

\bibitem{Surface}
L.~Eriksson,
Global ratio monotonicity for a killed von Mises bridge,
programme manuscript, July 2026;
\href{https://ai.vixra.org/abs/2607.0089}{ai.vixra:2607.0089}.

\end{thebibliography}

\end{document}
